\chapter{Marco teórico}

\section{Dinámica molecular}

La dinamica molecular (MD, por sus síglas en inglés) es una técnica de simulación por computadora que ha sido utilizada para analizar el comportamiento y evolución de diversos sistemas a nivel atomístico. Su esencia consiste en resolver numéricamente las ecuaciones clásicas del movimiento, típicamente las de Newton, para un conjunto de átomos y moléculas.Este método calcula las fuerzas, velocidades y posiciones de las partículas a lo largo de determinado tiempo (trayecotria) a partir de la cuál es posible derivar propiedades de interés mediante su análisis o experimentos computacionales posteriores.

\subsection{Condiciones iniciales}

Las condiciones iniciales en una simulación de dinámica molecular (MD) son el conjunto de posiciones y velocidades de todos los átomos del sistema en el instante de tiempo cero 